\documentclass[12pt]{article}
\usepackage[utf8]{inputenc}
\usepackage[T1]{fontenc}
\usepackage[margin=1in]{geometry}
\usepackage{setspace}
\usepackage{xcolor}
\usepackage[colorlinks=true, urlcolor=blue!60!black, linkcolor=blue!60!black]{hyperref}
\setlength{\parindent}{0em}
\setlength{\parskip}{0.6em}
\renewcommand{\baselinestretch}{1.0}
\begin{document}
\begin{center}
{\LARGE \textit{Tinker v. Des Moines Independent Community School District}, 393 U.S. 503 (1969)\par}
\vspace{0.5em}
{\large Case Brief --- Two-Page Edition\par}
\end{center}

\section{Memory Jogger}
Students keep their free-speech rights at school; passive political expression may be punished only on a reasonable forecast of material and substantial disruption (\href{https://supreme.justia.com/cases/federal/us/393/503/}{393 U.S. 503}).

\section{Detailed Case Facts}
In December 1965, Mary Beth Tinker (13), John Tinker (15), and Christopher Eckhardt (16) planned to wear black armbands to protest the Vietnam War. Des Moines principals preemptively banned armbands, though students wore campaign buttons and Iron Crosses freely. The three wore armbands, refused to remove them, and were suspended. Only a handful of 18,000 students took part, and there was no disruption of schoolwork.

\section{Procedural History}
The district court dismissed their \S~1983 suit, upholding the ban as reasonable; the en banc Eighth Circuit divided equally and affirmed; the Supreme Court reversed.

\section{Judicial Votes}
7--2. Fortas wrote for Warren, Douglas, Brennan, Stewart, White, and Marshall; Stewart and White concurred; Black and Harlan dissented.

\section{Holding}
Schools may not suppress passive student political expression without facts reasonably forecasting material and substantial interference with schoolwork or discipline, or invasion of others' rights; undifferentiated fear is not enough.

\section{Analysis of Opinions}
Fortas treated the silent armbands as close to pure speech, adopted the \textit{Burnside} standard, and found no evidence of forecast disruption, only a wish to avoid controversy, made worse by singling out antiwar symbols. Schools may not be enclaves of totalitarianism. Stewart doubted that children's rights equal adults'; White distinguished words from conduct. Black accused the Court of handing schools to their ``loudest-mouthed, but maybe not their brightest, students''; Harlan would have required proof of bad faith. \textit{Fraser}, \textit{Hazelwood}, \textit{Morse}, and \textit{Mahanoy} later refined the rule.

\section{Examples: Future Applications}
\textbf{Same-Side:} a quiet political T-shirt banned while opposing shirts are allowed is protected. \textbf{Opposite-Side:} a viewpoint-neutral apparel ban after documented fights is valid. \textbf{Fence-Sitter:} a weekend social-media post mocking a teacher that circulates in class.

\section{Critique}
Critics warn that the forecast-of-disruption test invites a heckler's veto and that the undefined rights-of-others prong divides courts.

\section{Key Quotations}
\begin{quote}
``It can hardly be argued that either students or teachers shed their constitutional rights to freedom of speech or expression at the schoolhouse gate.''
\end{quote}
\noindent --- Justice Fortas, Opinion of the Court, 393 U.S. at 506.

\end{document}
